\documentclass[10pt,a4paper]{amsart}
\usepackage[margin=19mm]{geometry}
\usepackage{amsmath,amssymb,mathtools}
\usepackage[scr=rsfs,cal=euler]{mathalfa}
\usepackage{xfrac}
\usepackage[unicode,hidelinks,bookmarks=false]{hyperref}
\newcommand{\ie}{i.e.\ }
\newcommand{\cf}{cf.\ }
\newcommand{\resp}{respectively\ }
\newcommand{\Subsection}{\subsection}
\providecommand{\nobreakdash}{\nobreak\hbox{-}}
\setlength{\emergencystretch}{2em}
\multlinegap0pt
\usepackage{latexsym}
\newtheorem{theorem}{Theorem}[section]
\newtheorem{lemma}{Lemma}[section]
\newtheorem{proposition}{Proposition}[section]
\newtheorem{corollary}{Corollary}[section]
\newtheorem{definition}{Definition}[section]
\newtheorem{remark}{Remark}[section]
\newtheorem{example}{Example}[section]
\numberwithin{equation}{section}
\newcommand{\R}{{\mathbb R}}
\newcommand{\N}{{\mathbb N}}
\newcommand{\D}{{\mathcal D}}
\newcommand{\const}{\mathrm{const}}
\newcommand{\supp}{\mathop{\rm supp}}
\newcommand{\codim}{\mathop{\rm codim}}
\renewcommand{\div}{\mathop{\rm div}\nolimits}
\renewcommand{\Im}{\mathop{\rm Im}}
\begin{document}
\title[On symmetric systems of transport equations]{On symmetric systems of transport equations}
\author{Evgeny Yu. Panov; of Mathematics, St. Petersburg, Russia; Veliky Novgorod, Russia}
\date{}
\maketitle
\noindent\emph{Source and licence.} On symmetric systems of transport equations. Version arXiv:2608.19835v1 (CC BY 4.0). This material is distributed under the Creative Commons Attribution 4.0 International licence, http://creativecommons.org/licenses/by/4.0/, as recorded in the arXiv OAI licence metadata for this exact version and shown on the abstract page https://arxiv.org/abs/2608.19835v1. Full rights notice: licenses/arxiv-2608.19835-CC-BY-4.0.txt.\par\smallskip
\noindent\textit{Editorial scope.} Counted unit: the original \texttt{theorem} environment \texttt{thM} at source lines 222--224 together with its complete proof at lines 226--232; the delivered document keeps source lines 92--233 so that every referenced label and every citation resolves inside it. Native editable source is the paper's own concatenated TeX; the AMS wrapper is editorial formatting only. No publisher class, upstream script or unrelated artwork is redistributed.
\begin{equation}\label{gr}
\|A^k(x)\|\le c(1+|x|), \quad k=1,\ldots,n, \ c=\const,
\end{equation}
where $\|A\|$ is the operator norm of a matrix $A$, and $|x|$ denotes the Euclidean norm of a finite-dimensional vector $x$.

By the definition of adjoint operator, $D(S^*)$ consists of such $u\in H$ that for all $f\in D(S_0)=C_0^1(\R^n,\R^m)$ the relation $(u,Sf)_H=(v,f)_H$ holds for some $v\doteq S^* u\in H$. Revealing the relation $(u,Sf)_H=(v,f)_H$, we obtain that
for all $f\in C_0^1(\R^n,\R^m)$
\begin{equation}\label{3}
\int_{\R^n} \sum_{k=1}^n u(x)\cdot A^k(x)f_{x_k}(x)dx=\int_{\R^n} v(x)\cdot f(x)dx.
\end{equation}
Since $u(x)\cdot A^k(x)f_{x_k}(x)=A^k(x)u(x)\cdot f_{x_k}(x)$ by symmetricity of the matrices $A^k$, relation (\ref{3}) means that
\begin{equation}\label{S*}
S^*u=v=-\sum_{k=1}^n (A^k(x)u(x))_{x_k} \ \mbox{ in } \D'(\R^n).
\end{equation}
Hence, the domain $D(S^*)$ consists of such $u(x)\in H$ that the distribution $\displaystyle v=-\sum_{k=1}^n (A^k(x)u(x))_{x_k}$ lies in $H$, and we set the value $S^*u=v$ as indicated in (\ref{S*}).

We denote by $H_c$ the space of vector-functions in $H$ with compact support.

\begin{lemma}\label{lem1}
The space $H_c\cap D(S^*)$ is a core of $S^*$, that is, it is dense in the space $D(S^*)$ equipped with the graph norm of $S^*$.
\end{lemma}

\begin{proof}
Assume that $u(x)\in D(S^*)$. We need to approximate $u(x)$ in the graph norm of $D(S^*)$ by compactly supported vector-functions in $D(S^*)$. Let $p(y)\in C_0^1(\R^n)$ be such a function that $p(y)=1$ for $|y|\le 1$, $p(y)=0$ for $|y|\ge 2$. We set $u_R(x)=p(x/R)u(x)$, where $R>1$. It is clear that $u_R\in H_c$.
Since $p(x/R)\in C_0^1(\R^n)$, the distribution $\displaystyle v_R\doteq -\sum_{k=1}^n (A^k(x)p(x/R)u(x))_{x_k}$ can be represented as
\begin{equation}\label{4}
v_R=-p(x/R)\sum_{k=1}^n (A^k(x)u(x))_{x_k}-\frac{1}{R}\sum_{k=1}^n p_{y_k}(x/R)A^k(x)u(x)=p(x/R)v(x)-w_R(x),
\end{equation}
where $v(x)=S^*u(x)\in H$, $\displaystyle w_R(x)=\frac{1}{R}\sum_{k=1}^n p_{y_k}(x/R)A^k(x)u(x)\in H$. We see that $v_R\in H$ and therefore $u_R\in D(S^*)$, $S^*u_R=v_R$ (see the description of $D(S^*)$ above). Thus, $u_R\in H_c\cap D(S^*)$. It is clear that $u_R=p(x/R)u(x)\to u(x)$, $p(x/R)v(x)\to v(x)$ in $H$ as $R\to\infty$. Let us show that $w_R(x)\to 0$ in $H$
as $R\to\infty$. Notice that $\supp \nabla_y p(x/R)$ lie in the annulus $R\le |x|\le 2R$ while $\|A^k(x)\|\le c(1+2R)\le 3cR$ in this annulus, by condition (\ref{gr}). Therefore,
\[
|w_R(x)|\le 3cn\|\nabla p\|_\infty |u(x)|\theta(|x|-R),
\]
where $\theta(s)$ is the Heaviside function. This estimate implies the desired relation $w_R\to 0$ in $H$ as $R\to\infty$.
It follows that $v_R\mathop{\to}\limits_{R\to\infty} v$ in $H$. Hence, $u_R\to u$, $S^*u_R\to S^*u$ in $H$, which means
convergence $u_R\to u$ in the graph norm. The proof is complete.
\end{proof}

We need a variant of the famous DiPerna-Lions commutation lemma
\cite[Lemma~II.1]{DiL}. Suppose $\rho(z)\in C_0^1(\R^n)$, $\supp\rho\subset B_1(0)=\{z\in\R^n | |z|\le 1\}$, $\rho(z)\ge 0$, $\int_{\R^n}\rho(z)dz=1$. For $u(x)\in L^1_{loc}(\R^n)$ we introduce the corresponding averaged functions
\[u_h(x)=u*\rho_h(x)=\int_{\R^n}\rho_h(x-y)u(y)dy,\]
where $h\in (0,1)$, $\rho_h(z)=h^{-n}\rho(z/h)$.

\begin{lemma}\label{lmDL}
Let $a(x)\in C(\R^n)$ satisfy the Lipschitz condition
${|a(x)-a(y)|\le L|x-y|}$ for all $x,y\in\R^n$; $u(x)\in L^p(\R^n)$, $1\le p<\infty$. Then
\begin{equation}\label{DL}
\frac{\partial}{\partial x_k}(au_h-(au)_h)(x)\mathop{\to}_{h\to 0} 0 \ \mbox{ in } L^p(\R^n)
\end{equation}
for all $k=1,\ldots,n$.
\end{lemma}

\begin{proof}
First of all we notice that by the Lipschitz condition the generalized derivatives $a_{x_k}(x)\in L^\infty(\R^n)$, $k=1,\dots,n$, and $\|\nabla a(x)\|_\infty\le L$. Since \[(au_h-(au)_h)(x)=h^{-n}\int_{\R^n}(a(x)-a(y))\rho((x-y)/h) u(y)dy,\]
there exist the generalized derivatives
\begin{align}\label{rel}
\frac{\partial}{\partial x_k}(au_h-(au)_h)(x)=h^{-n}\int_{\R^n}a_{x_k}(x)\rho((x-y)/h) u(y)dy+\nonumber\\ h^{-n-1}\int_{\R^n}(a(x)-a(y))\rho_{z_k}((x-y)/h) u(y)dy=I_1^h(x)+I_2^h(x).
\end{align}
The first term in this sum
\begin{equation}\label{rel1}
I_1^h(x)=a_{x_k}(x)\int_{\R^n}\rho_h(x-y) u(y)dy=a_{x_k}(x)u_h(x)\to a_{x_k}(x)u(x)
\end{equation}
as $h\to 0$ in $L^p(\R^n)$ because $u_h\mathop{\to}\limits_{h\to 0} u$ in $L^p(\R^n)$ by the known property of averaged functions while the derivative $a_{x_k}(x)\in L^\infty(\R^n)$.
Next, we estimate the term $I_2^h(x)$. Obviously,
\begin{align}\label{es1}
|I_2^h(x)|\le h^{-n-1}\int_{\R^n}|a(x)-a(y)||\rho_{z_k}((x-y)/h)| |u(y)|dy\le \nonumber\\ \omega^h(x)\doteq Lh^{-n}\int_{\R^n} h^{-1}|x-y||\rho_{z_k}(h^{-1}(x-y))| |u(y)|dy.
\end{align}
By the properties of averaged functions (with the kernel $|z||\rho_{z_k}(z)|$) $\omega^h(x)\in L^p(\R^n)$ and
\begin{equation}\label{rel2}
\omega^h(x)\to C|u(x)|
\end{equation}
as $h\to 0$ both in $L^p(\R^n)$ and a.e. in $\R^n$. Here $\displaystyle C=L\int_{\R^n} |z||\rho_{z_k}(z)|dz=\const$.
Now, let $x$ be a common Lebesgue point of the vector $\nabla a(y)$ and the function $u(y)$. Then
\begin{align}\label{rel2-1}
I_2^h(x)=h^{-n-1}\int_{\R^n}(a(x)-a(y))\rho_{z_k}(h^{-1}(x-y))(u(y)-u(x))dy+\nonumber\\ u(x)h^{-n-1}\int_{\R^n}(a(x)-a(y))\rho_{z_k}(h^{-1}(x-y))dy.
\end{align}
The first term in the right-hand part of (\ref{rel2-1}) is estimated as
\begin{align}\label{est2-1}
h^{-n-1}\left|\int_{\R^n}(a(x)-a(y))\rho_{z_k}(h^{-1}(x-y))(u(y)-u(x))dy\right|\le \nonumber\\
Lh^{-n}\int_{\R^n}h^{-1}|x-y||\rho_{z_k}(h^{-1}(x-y))||u(y)-u(x)|dy\mathop{\to}_{h\to 0} 0
\end{align}
since $x$ is a Lebesgue point of $u(y)$. We introduce the function ${\displaystyle J^h(x)=h^{-n-1}\int_{\R^n}(a(x)-a(y))\rho_{z_k}(h^{-1}(x-y))dy}$ and represent it in the form
\begin{align}\label{rel2-2}
J^h(x)=h^{-n-1}\int_{\R^n}(a(x)-a(y)-\nabla a(x)\cdot(x-y))\rho_{z_k}(h^{-1}(x-y))dy+\nonumber\\ h^{-n-1}\int_{\R^n}\nabla a(x)\cdot (x-y)\rho_{z_k}(h^{-1}(x-y))dy=\nonumber\\
h^{-n-1}\int_{\R^n}\int_0^1 (\nabla a(x+s(y-x))-\nabla a(x))\cdot(x-y)\rho_{z_k}(h^{-1}(x-y))dsdy+\nonumber\\ \nabla a(x)\cdot\int_{\R^n}z\rho_{z_k}(z)dz.
\end{align}
We utilize here the identity
\[
a(x)-a(y)=\int_0^1\nabla a(x+s(y-x))\cdot(x-y)ds,
\]
which holds for a.e. $y\in\R^n$.
To estimate the first term in the right-hand part of (\ref{rel2-2}), we make the change of variables $y\to z=s(x-y)$, resulting in
\begin{align}\label{est2-2}
h^{-n-1}\left|\int_{\R^n}\int_0^1 (\nabla a(x+s(y-x))-\nabla a(x))\cdot(x-y)\rho_{z_k}(h^{-1}(x-y))dsdy\right|\le\nonumber\\
h^{-n}\int_{\R^n} |\nabla a(x-z)-\nabla a(x)|\rho_1(z/h)dz,
\end{align}
where we denote
\[
\rho_1(y)=|y|\int_0^1s^{-n-1}|\rho_{z_k}(y/s)|ds.
\]
Remark that $\rho_1\ge 0$, $\supp\rho_1\subset B_1(0)$, 
\begin{align*}
\int_{\R^n}\rho_1(y)dy=\int_0^1\int_{\R^n} |s^{-1}y||\rho_{z_k}(s^{-1}y)|s^{-n}dyds= \\ \int_0^1\int_{\R^n}|z||\rho_{z_k}(z)|dz ds=\int_{\R^n}|z||\rho_{z_k}(z)|dz
\end{align*}
(we made the change $z=s^{-1}y$), and since $x$ is a Lebesgue point of the vector $\nabla a(y)$, we find
\[h^{-n}\int_{\R^n} |\nabla a(x-z)-\nabla a(x)|\rho_1(z/h)dz\mathop{\to}_{h\to 0} 0.\]
It follows from (\ref{rel2-2}) and (\ref{est2-2}) that
\begin{equation} \label{rel3}
J^h(x)\mathop{\to}_{h\to 0} \nabla a(x)\cdot\int_{\R^n}z\rho_{z_k}(z)dz=-a_{x_k}(x).
\end{equation}
Here, we apply the integration by parts formula
\[\int_{\R^n}z\rho_{z_k}(z)dz=-\int_{\R^n}\frac{\partial z}{\partial z_k}\rho(z)dz=-e_k, \]
where $e_k$ is the $k$th basis vector in $\R^n$.
In view of (\ref{rel2-1}), (\ref{est2-1}) it follows from (\ref{rel3}) that $I_2^h(x)\mathop{\to}\limits_{h\to 0} -a_{x_k}(x)u(x)$. By our choice, $x\in\R^n$ is an arbitrary point of a set of full Lebesgue measure. Therefore, $I_2^h(x)\to -a_{x_k}(x)u(x)$ as $h\to 0$ a.e. in $\R^n$. Further, we notice that
\begin{align*}
|I_2^h(x)+a_{x_k}(x)u(x)|^p\le 2^{p-1}(|I_2^h(x)|^p+|a_{x_k}(x)u(x)|^p)\le \\ 2^{p-1}((\omega^h(x))^p+|a_{x_k}(x)u(x)|^p).
\end{align*}
The left-hand side of this inequality converges as $h\to 0$ to zero a.e. in $\R^n$, while its right-hand side converges both in  $L^1(\R^n)$ and a.e. in $\R^n$, by relation (\ref{rel2}). Applying Fatou's lemma to the sequence
\[
2^{p-1}((\omega^h(x))^p+|a_{x_k}(x)u(x)|^p)-|I_2^h(x)+a_{x_k}(x)u(x)|^p,
\]
we derive that
\[
\lim_{h\to 0} \int_{\R^n}|I_2^h(x)+a_{x_k}(x)u(x)|^pdx=0,
\]
that is, $\displaystyle I_2^h(x)\mathop{\to}_{h\to 0} -a_{x_k}(x)u(x)$ in $L^p(\R^n)$. This, together with  (\ref{rel}), (\ref{rel1}), gives the required relation (\ref{DL}). The proof is complete.
\end{proof}

\section{Main results}\label{secM}


\begin{theorem}\label{thM}
Assume that the coefficients $a^k_{ij}(x)\in W^1_{\infty,loc}(\R^n)$ and satisfy the growth condition $(1+|x|)^{-1}a^k_{ij}(x)\in L^\infty(\R^n)$. Then the operator $S$ is skew-adjoint.
\end{theorem}

\begin{proof}
We are going to demonstrate that the space $H_c\cap D(S^*)\subset D(S)$. Choosing $u=u(x)\in H_c\cap D(S^*)$, we consider the sequence of averaged vector-functions $u_h(x)=u*\rho_h(x)\in C_0^1(\R^n,\R^m)=D(S_0)$. If $\supp u(x)$ is contained in a ball $B_R(0)$ then $\supp u_h\subset B_{R+1}(0)$ for all $h\in (0,1)$ and the vectors $v=S^*u$, $v^h\doteq S^* u_h=-S_0u_h$ do not depend on matrices $A^k(x)$ for $|x|>R+1$. By our assumptions these matrices are Lipschitz continuous in $B_{R+1}(0)$ and can be extended to Lipschitz continuous symmetric matrices on the whole space $\R^n$. Hence, without loss of generality we may initially suppose that the matrices $A^k(x)$ satisfy the global Lipschitz condition in $\R^n$.
It then follows from Lemma~\ref{lmDL} that for all $i=1,\ldots,m$
\[(v^h-v_h)_i(x)=\sum_{k=1}^n \sum_{j=1}^m \frac{\partial}{\partial x_k}[(a^k_{ij}u_j)_h(x)-a^k_{ij}(x)(u_j)_h(x)]
\mathop{\to}_{h\to 0} 0 \ \mbox{ in } L^2(\R^n). \]
This implies that $v^h-v_h\to 0$ as $h\to 0$ in the space $H$. On the other hand, $u_h\mathop{\to}\limits_{h\to 0} u$,  $v_h\mathop{\to}\limits_{h\to 0} v$ in $H$ by the property of averaged functions. Hence, $u_h\to u$, $v^h=-Su_h\to v$ as $h\to 0$ in $H$. Since the operator $S$ is closed, we conclude that $u\in D(S)$ and $-Su=v$. We have proved the required inclusion $H_c\cap D(S^*)\subset D(S)$. Since the operator $S$ is closed, the closure of $H_c\cap D(S^*)$ in $D(S^*)$ with respect to the graph norm also lies in $D(S)$. But by Lemma~\ref{lem1} this closure coincides with $D(S^*)$, and we obtain that $D(S^*)\subset D(S)\subset D(S^*)$. Thus, $D(S)=D(S^*)$, therefore $-S=S^*$ and $S$ is a skew-adjoint operator. The proof is complete.
\end{proof}


\begin{thebibliography}{99}
\bibitem[DiL]{DiL} R.J.~DiPerna, P.L.~Lions, Ordinary differential equations, transport theory and Sobolev spaces, Invent. Math. 98 (1989) 511--547.
\end{thebibliography}
\end{document}
